\begin{abstract}
In healthcare, Electronic Health Records (EHRs) and continuous patient monitoring data are restricted by strict privacy regulations like GDPR in the European Union and HIPAA in the USA. To overcome this obstacle in medical AI development, we developed a rule-based synthetic data generation program that simulates a real-world Internet of Medical Things (IoMT) environment. The pipeline generated 500,000 synthetic patient records, including continuous vital signs, demographic history, and laboratory biomarkers across 24 distinct medical pathologies. To demonstrate the clinical representativeness of this dataset, we validated its statistical distributions against a real-world open-access Emergency Department admission dataset. Furthermore, we benchmarked an Early Fusion Multi-Layer Perceptron (MLP), a Contextual Fusion Tabular Transformer, and a tree-based baseline (XGBoost). Our results demonstrate that the synthetic data poses a challenging classification task. While models achieved a high ROC-AUC by effectively mapping multidimensional boundaries, deliberate overlapping clinical phenotypes mirrored real-world diagnostic ambiguity, limiting Macro F1-scores appropriately. Finally, SHapley Additive exPlanations (SHAP) confirmed that the algorithms learned genuine expert-driven physiological rules, suggesting this framework as a useful tool for privacy-preserving medical predictive modeling.
\end{abstract}
\begin{IEEEkeywords}
Synthetic Data Generation, Electronic Health Records (EHR), Internet of Medical Things (IoMT), Multimodal Data Fusion, Tabular Transformer, Explainable AI.
\end{IEEEkeywords}